\subsection{Formatting SQL inside Java}\label{sec:FormattingSQLInsideJava}
If your code needs to access a database\index{database} on an RDBMS\index{RDBMS}, it is quite likely that you have to write some SQL\index{SQL} \autocite{ANSI:SQLStandard,W3SCHOOLS:SQL} statements into your code. Even when you are using an ORM\index{ORM} like JPA\index{JPA} \autocite{ECLIPSE_JPA} and Hibernate\index{Hibernate} \autocite{HIBERNATE_ORM}, it might still be necessary to add some SQL\index{SQL} statements or statements in the tool's variant of SQL (HQL\index{HQL}\index{HQL|see {Hibernate}} in case of Hibernate\index{Hibernate}).

Also the statements of several other query languages can be formatted in the same manner than SQL\index{SQL} statements.

\keeplisting{6}
The SQL \verb#SELECT#\index{SQL!SELECT} statement
\begin{lstlisting}[language=SQL]
SELECT person.id AS "id", person.name AS "lastName", person.surname AS "firstName", address.street AS "street", address.city AS "city"
    FROM contact.person, contact.address
    WHERE person.addressid = address.id
        AND person.profession like '%Teacher%';
\end{lstlisting}

\keeplisting{8}
would be embedded into a Java source file like this:
\begin{lstlisting}
final var sqlSelect = """
````SELECT person.id AS "id", person.name AS "lastName", person.surname AS "firstName", address.street AS "street", address.city AS "city"
````    FROM contact.person, contact.address
````    WHERE person.addressid = address.id
````        AND person.profession like '%Teacher%';\
````""";
\end{lstlisting}

Basically, the keywords are written all uppercase (\verb#SELECT#\index{SQL!SELECT}, \verb#FROM#\index{SQL!FROM},~…), while the arguments are all lowercase (\verb#person.name#, \verb#contact.address#,~…). Alias\index{SQL!alias} names are \textit{always} surrounded by double quotes even they do not contain blanks or are SQL\index{SQL} keywords\footnote{Some RDBMS\index{RDBMS} products do not allow double quotes for this purpose, they require square brackets.}.

The indendation\index{indentation} (four spaces)\index{indendation!character}\index{indentation!size} is per level: \verb#FROM#\index{SQL!FROM} and \verb#WHERE#\index{SQL!WHERE} are one level down from \verb#SELECT#\index{SQL!SELECT}, \verb#AND#\index{SQL!AND} is one level down from \verb#WHERE#\index{SQL!WHERE}.

%==============================================================================
% End of file!!
%==============================================================================
